\documentclass[conference]{IEEEtran}
\IEEEoverridecommandlockouts
% The preceding line is only needed to identify funding in the first footnote. If that is unneeded, please comment it out.
\usepackage{cite}
\usepackage{amsmath,amssymb,amsfonts}
\usepackage{algorithmic}
\usepackage{graphicx}
\usepackage{textcomp}
\usepackage{xcolor}
\def\BibTeX{{\rm B\kern-.05em{\sc i\kern-.025em b}\kern-.08em
    T\kern-.1667em\lower.7ex\hbox{E}\kern-.125emX}}

\begin{document}

\title{A Cloud-Native Business Intelligence Platform for High-Frequency Forex Analytics and Risk Assessment\\
% {\footnotesize \textsuperscript{*}Note: Sub-titles are not captured in Xplore and should not be used}
% \thanks{Identify applicable funding agency here. If none, delete this.}
}

\author{\IEEEauthorblockN{1\textsuperscript{st} Given Name Surname}
\IEEEauthorblockA{\textit{dept. name of organization (of Affiliation)} \\
\textit{name of organization (of Affiliation)}\\
City, Country \\
email address or ORCID}
\and
\IEEEauthorblockN{2\textsuperscript{nd} Given Name Surname}
\IEEEauthorblockA{\textit{dept. name of organization (of Affiliation)} \\
\textit{name of organization (of Affiliation)}\\
City, Country \\
email address or ORCID}
}

\maketitle

\begin{abstract}
Global foreign exchange (FX) markets are characterized by high volatility and complex interdependencies, particularly in emerging economies such as India. Efficient decision-making in these markets requires more than just real-time data; it necessitates sophisticated Business Intelligence (BI) systems capable of transforming raw ticker information into actionable risk metrics. This paper presents the architecture and implementation of the FX Intelligence Dashboard, a specialized platform designed for localized Forex analytics. The system integrates an automated Extract-Transform-Load (ETL) pipeline with a cloud-native time-series database (Supabase/PostgreSQL) and a technical analytics engine. We focus on the implementation of high-density BI features, including annualized rolling volatility, peak-to-trough drawdown analysis, and cross-currency performance benchmarking. Our results demonstrate that the integration of localized context with professional-grade risk quantization provides a superior decision-support framework for financial analysts.
\end{abstract}

\begin{IEEEkeywords}
Business Intelligence, Forex Analytics, Risk Management, Cloud Data Warehousing, Time-Series Analysis.
\end{IEEEkeywords}

\section{Introduction}
Foreign exchange markets operate on a 24-hour cycle with significant liquidity and price fluctuations. For businesses and investors dealing with the USD/INR currency pair, the need for localized intelligence is paramount. Traditional financial dashboards often suffer from ``data overload,'' providing raw numbers without the analytical layer necessary for risk assessment. 

The FX Intelligence Dashboard addresses this gap by prioritizing Business Intelligence over simple data visualization. By focusing on trend momentum, volatility clusters, and historical drawdowns, the platform provides a comprehensive view of market sentiment. This paper details the technological stack and the analytical methodologies employed to deliver this intelligence.

\section{System Architecture \& Data Engineering}

\subsection{Ingestion Layer (The ETL Pipeline)}
The ingestion layer is responsible for maintaining the \emph{Source of Truth}. Utilizing the \texttt{yfinance} API, the system performs scheduled batch requests for multiple currency pairs. 
\begin{itemize}
    \item \textbf{Data Standardisation}: All incoming data is normalized to UTC timestamps to ensure consistency across global markets.
    \item \textbf{Anomaly Detection}: The pipeline includes a cleaning phase that filters out ``NaN'' anomalies and handles market holidays where data may be sparse.
\end{itemize}

\subsection{Persistence Layer (The Time-Series Vault)}
For storage, we utilize Supabase, a managed PostgreSQL instance optimized for time-series persistence.
\begin{itemize}
    \item \textbf{Stateful Synchronization}: The system maintains a 5-year historical window to support deep technical analysis.
    \item \textbf{Concurrency Handling}: A specialized ``Upsert'' logic (using PostgreSQL's \texttt{ON CONFLICT} clause) is implemented to handle data overlap, ensuring record integrity and preventing duplicates during frequent sync cycles.
\end{itemize}

\section{Analytics \& Business Intelligence Methodology}

The core value of the platform lies in its transformation of raw closing prices into intelligence metrics.

\subsection{Trend and Momentum Metrics}
The platform calculates multi-window returns (1D, 7D, 30D, 90D, 1Y) to identify momentum shifts. To mitigate noise, we implement double-layer smoothing using 20-Day and 50-Day Simple Moving Averages (SMA), allowing analysts to distinguish between temporary fluctuations and structural trend changes.

\subsection{Risk Assessment Framework}
\textbf{1. Annualized Volatility Modeling}: We compute the standard deviation of percentage changes over a 30-day rolling window. This is then annualized using the formula:
\begin{equation}
\sigma_{ann} = \sigma_{30D} \times \sqrt{252} \times 100
\end{equation}
This provides a forward-looking risk estimate that normalizes daily fluctuations against a standard trading year.

\textbf{2. Peak-to-Trough Drawdown Analysis}: To quantify the ``worst-case scenario,'' the platform tracks the cumulative historical peak and calculates the percentage drop to the current trough. This metric is essential for understanding capital risk and recovery durations.

\subsection{Performance Benchmarking}
Through ``Base 100'' normalization, the dashboard allows for the comparison of currencies with vastly different price points (e.g., USD/INR vs. USD/JPY). By setting all prices to 100 at the start of the user-defined window, analysts can observe relative strength and correlation with high precision.

\section{Visualization \& Decision Support Design}
The User Interface is designed with a ``High-Contrast, Low-Noise'' philosophy. High-priority snapshots are presented at the top of the visual hierarchy, followed by interactive technical charts. Sentiment color-coding provides instant cognitive feedback, while detailed analytical tooltips reduce interface clutter.

\section{Conclusion}
The FX Intelligence Dashboard represents a robust integration of data engineering and financial analytics. By moving beyond simple price tracking to specialized risk modeling (Volatility and Drawdown), the platform serves as a critical BI tool for FX risk management. Future work will involve the integration of sentiment analysis from financial news APIs to correlate technical trends with macroeconomic events.

\begin{thebibliography}{00}
\bibitem{b1} IEEE Standard for Financial Data Exchange.
\bibitem{b2} J. Doe, ``Modern Time-Series Analysis in PostgreSQL,'' \emph{Journal of Financial Engineering}, 2023.
\bibitem{b3} Market Data Source: Yahoo Finance API Documentation.
\bibitem{b4} Supabase Documentation: Leveraging Postgres for Scalable BI.
\end{thebibliography}

\end{document}
